\scanpage{233}
est l'hom. composé [linéaire pour la structure \uline{gauche} sur $P^i(M)$]
\diagram{\begin{tikzcd}
E\arrow[r,"\delta_i"]&P^i(E)\arrow[r,no head]&P^i(M).
\end{tikzcd}}
Donc, pour $P$ variable, c'est représenté par
\[
\Phi_i=P^i(M)/\underline{O}_S.\rho_i(E),
\]
les $\Phi$ forment un syst. projectif strict, déterminé par la connaissance de $M$ et de
\[
\rho_\infty:E\longrightarrow P^\infty(M)
\]
(linéaire pour la structure \uline{gauche} de $P^\infty(M)$), défini par les $\rho_i$.
Attention, ici $\rho_\infty$ satisfait à des \deleted{\unsure{le}} conditions spéciales: \unsure{son} composé $E\to M$ avec l'augmentation $P^\infty(M)\to M$ est un épimorphisme, $+$ conditions d'intégrabilité ---
\margin{232}
